\documentclass[10pt,a4paper]{amsart}
\usepackage{amsmath,amssymb,amsthm,graphicx}
% Retained preamble for the Stacks Project chapter {CH} (shared project preamble).
% The chapter file pinned-original-source.tex only carries \input{{preamble}} on its first line,
% so no preamble text lives inside the chapter itself.  The Stacks Project uses one single shared
% preamble for every chapter: the snapshot stages it three times (sources/stacks-chow/preamble.tex,
% sources/stacks-etale-cohomology/preamble.tex and focused-seeds/AGP/sources/stacks-snapshot-20261011/
% preamble.tex) and those copies are identical in size.  The block below is copied mechanically,
% line for line, from sources/stacks-chow/preamble.tex lines 14-18 and 154-217: the Xy-pic setup and
% the native theorem declarations together with the mathematical macros.  Nothing was added, removed
% or redefined; the wrapper class and the amsmath/amssymb/amsthm/graphicx base packages are supplied
% by the assembler.
% For commutative diagrams we use Xy-pic
% Chapter: stacks-spaces-more-morphisms.tex -- More on Morphisms of Spaces.
\usepackage[all]{xy}

% We use 2cell for 2-commutative diagrams.
\xyoption{2cell}

% Theorem environments.
%
\theoremstyle{plain}
\newtheorem{theorem}[subsection]{Theorem}
\newtheorem{proposition}[subsection]{Proposition}
\newtheorem{lemma}[subsection]{Lemma}

\theoremstyle{definition}
\newtheorem{definition}[subsection]{Definition}
\newtheorem{example}[subsection]{Example}
\newtheorem{exercise}[subsection]{Exercise}
\newtheorem{situation}[subsection]{Situation}

\theoremstyle{remark}
\newtheorem{remark}[subsection]{Remark}
\newtheorem{remarks}[subsection]{Remarks}

\numberwithin{equation}{subsection}

% Macros
%
\def\lim{\mathop{\mathrm{lim}}\nolimits}
\def\colim{\mathop{\mathrm{colim}}\nolimits}
\def\Spec{\mathop{\mathrm{Spec}}}
\def\Hom{\mathop{\mathrm{Hom}}\nolimits}
\def\Ext{\mathop{\mathrm{Ext}}\nolimits}
\def\SheafHom{\mathop{\mathcal{H}\!\mathit{om}}\nolimits}
\def\SheafExt{\mathop{\mathcal{E}\!\mathit{xt}}\nolimits}
\def\Sch{\mathit{Sch}}
\def\Mor{\mathop{\mathrm{Mor}}\nolimits}
\def\Ob{\mathop{\mathrm{Ob}}\nolimits}
\def\Sh{\mathop{\mathit{Sh}}\nolimits}
\def\NL{\mathop{N\!L}\nolimits}
\def\CH{\mathop{\mathrm{CH}}\nolimits}
\def\proetale{{pro\text{-}\acute{e}tale}}
\def\etale{{\acute{e}tale}}
\def\QCoh{\mathit{QCoh}}
\def\Ker{\mathop{\mathrm{Ker}}}
\def\Im{\mathop{\mathrm{Im}}}
\def\Coker{\mathop{\mathrm{Coker}}}
\def\Coim{\mathop{\mathrm{Coim}}}

% Boxtimes
%
\DeclareMathSymbol{\boxtimes}{\mathbin}{AMSa}{"02}

%
% Macros for moduli stacks/spaces
%
\def\QCohstack{\mathcal{QC}\!\mathit{oh}}
\def\Cohstack{\mathcal{C}\!\mathit{oh}}
\def\Spacesstack{\mathcal{S}\!\mathit{paces}}
\def\Quotfunctor{\mathrm{Quot}}
\def\Hilbfunctor{\mathrm{Hilb}}
\def\Curvesstack{\mathcal{C}\!\mathit{urves}}
\def\Polarizedstack{\mathcal{P}\!\mathit{olarized}}
\def\Complexesstack{\mathcal{C}\!\mathit{omplexes}}
% \Pic is the operator that assigns to X its picard group, usage \Pic(X)
% \Picardstack_{X/B} denotes the Picard stack of X over B
% \Picardfunctor_{X/B} denotes the Picard functor of X over B
\def\Pic{\mathop{\mathrm{Pic}}\nolimits}
\def\Picardstack{\mathcal{P}\!\mathit{ic}}
\def\Picardfunctor{\mathrm{Pic}}
\def\Deformationcategory{\mathcal{D}\!\mathit{ef}}
\newenvironment{reference}{}{}
\newenvironment{slogan}{}{}
\newenvironment{history}{}{}
\usepackage{cleveref}
\crefname{originallemma}{original lemma}{original lemmas}
\Crefname{originallemma}{Original lemma}{Original lemmas}
\makeatletter
\newlabel{lemma-chow-noetherian}{{088S}{source}{}{sourcepointer-1}{}}
\newlabel{lemma-chow-noetherian@cref}{{[originallemma][0][]088S}{[1][0][]source}{}{}{}}
\makeatother
\crefname{originallemma}{original lemma}{original lemmas}
\Crefname{originallemma}{Original lemma}{Original lemmas}
\makeatletter
\newlabel{morphisms-lemma-H-projective-composition}{{01WE}{source}{}{sourcepointer-2}{}}
\newlabel{morphisms-lemma-H-projective-composition@cref}{{[originallemma][0][]01WE}{[1][0][]source}{}{}{}}
\makeatother
\crefname{originallemma}{original lemma}{original lemmas}
\Crefname{originallemma}{Original lemma}{Original lemmas}
\makeatletter
\newlabel{morphisms-lemma-projective-over-quasi-projective-is-H-projective}{{087S}{source}{}{sourcepointer-3}{}}
\newlabel{morphisms-lemma-projective-over-quasi-projective-is-H-projective@cref}{{[originallemma][0][]087S}{[1][0][]source}{}{}{}}
\makeatother
\crefname{originallemma}{original lemma}{original lemmas}
\Crefname{originallemma}{Original lemma}{Original lemmas}
\makeatletter
\newlabel{spaces-limits-lemma-descend-finite-presentation}{{07SK}{source}{}{sourcepointer-4}{}}
\newlabel{spaces-limits-lemma-descend-finite-presentation@cref}{{[originallemma][0][]07SK}{[1][0][]source}{}{}{}}
\makeatother
\crefname{originallemma}{original lemma}{original lemmas}
\Crefname{originallemma}{Original lemma}{Original lemmas}
\makeatletter
\newlabel{spaces-limits-lemma-descend-separated}{{084T}{source}{}{sourcepointer-5}{}}
\newlabel{spaces-limits-lemma-descend-separated@cref}{{[originallemma][0][]084T}{[1][0][]source}{}{}{}}
\makeatother
\crefname{originallemma}{original lemma}{original lemmas}
\Crefname{originallemma}{Original lemma}{Original lemmas}
\makeatletter
\newlabel{spaces-limits-lemma-descend-separated-morphism}{{0851}{source}{}{sourcepointer-6}{}}
\newlabel{spaces-limits-lemma-descend-separated-morphism@cref}{{[originallemma][0][]0851}{[1][0][]source}{}{}{}}
\makeatother
\crefname{originallemma}{original lemma}{original lemmas}
\Crefname{originallemma}{Original lemma}{Original lemmas}
\makeatletter
\newlabel{spaces-limits-lemma-topology-limit}{{086V}{source}{}{sourcepointer-7}{}}
\newlabel{spaces-limits-lemma-topology-limit@cref}{{[originallemma][0][]086V}{[1][0][]source}{}{}{}}
\makeatother
\crefname{originallemma}{original lemma}{original lemmas}
\Crefname{originallemma}{Original lemma}{Original lemmas}
\makeatletter
\newlabel{spaces-limits-proposition-approximate}{{07SU}{source}{}{sourcepointer-8}{}}
\newlabel{spaces-limits-proposition-approximate@cref}{{[originallemma][0][]07SU}{[1][0][]source}{}{}{}}
\makeatother
\crefname{originallemma}{original lemma}{original lemmas}
\Crefname{originallemma}{Original lemma}{Original lemmas}
\makeatletter
\newlabel{spaces-morphisms-lemma-quasi-compact-scheme-theoretic-image}{{082Z}{source}{}{sourcepointer-9}{}}
\newlabel{spaces-morphisms-lemma-quasi-compact-scheme-theoretic-image@cref}{{[originallemma][0][]082Z}{[1][0][]source}{}{}{}}
\makeatother
\title{Chow's lemma for algebraic spaces: a blowup of a separated finite type space embeds into projective space over the base}
\author{The Stacks Project Authors}
\date{stacks-project:spaces-more-morphisms.tex:local-snapshot-20261011; GNU FDL 1.2 or later}
\begin{document}
\maketitle
\noindent\textit{Editorial note.} Statement span is lines 9274-9290 of sources/stacks-spaces-more-morphisms/pinned-original-source.tex; proof span is lines 9292-9365, the complete author proof (72 body lines).\par\medskip
\noindent\textit{Editorial note.} Navigation stubs: 9 cited-but-not-included labels are supplied as non-load-bearing pointers with their pinned tags where the snapshot tag index has them.\par\medskip
\noindent\textit{Editorial note.} extra\_packages is empty on purpose: the retained preamble already loads xy with options and a second declaration clashes in the assembler wrapper.\par\medskip
\noindent\textit{Editorial note.} Text is retained verbatim, including the author's own typos and placeholders.\par\medskip
\noindent\textit{Editorial note.} The chapter cites Knutson, IV Theorem 3.1 in a reference block; that citation is part of the statement block and no bibliography entry is required for the collected spans (no \textbackslash{}cite outside the reference environment).\par\medskip
\noindent\textit{Editorial note.} The statement's reference environment cites \textbackslash{}cite[IV Theorem 3.1]\{Kn\}; a bibliography.tex holding the snapshot my.bib Kn entry (Knutson, Algebraic spaces, LNM 203, Springer-Verlag 1971) is attached (precedent S3-morph-etale-nbhd-dominates-smooth EGA4).\par\medskip
\noindent\textit{Allowed standard prerequisites (editorial).}\begin{itemize}
\item Explicit prerequisite boundary. Approximation and descent of spaces and morphisms of finite presentation (Limits of Spaces, proposition-approximate, lemmas descend-separated, descend-finite-presentation, descend-separated-morphism, topology-limit), the schemes-level Chow input and the H-projectivity criteria from the Morphisms chapter (projective-over-quasi-projective-is-H-projective, H-projective-composition); cited inputs with the reduction and diagram construction written in the retained span.
\end{itemize}\medskip
\ifdefined\hypertarget\hypertarget{sourcepointer-1}{}\fi
\noindent\textit{Source reference.} Original lemma 088S: sources/stacks-spaces-more-morphisms/pinned-original-source.tex, source lines 9225--9244. Cited by the retained statement or proof but not part of the collected spans; supplied only as a navigation pointer. The argument does not depend on its content.\par\medskip
\ifdefined\hypertarget\hypertarget{sourcepointer-2}{}\fi
\noindent\textit{Source reference.} Original lemma 01WE: sources/stacks-spaces-more-morphisms/pinned-original-source.tex, source lines 9365--9365. Cited by the retained statement or proof but not part of the collected spans; supplied only as a navigation pointer. The argument does not depend on its content.\par\medskip
\ifdefined\hypertarget\hypertarget{sourcepointer-3}{}\fi
\noindent\textit{Source reference.} Original lemma 087S: sources/stacks-spaces-more-morphisms/pinned-original-source.tex, source lines 9365--9365. Cited by the retained statement or proof but not part of the collected spans; supplied only as a navigation pointer. The argument does not depend on its content.\par\medskip
\ifdefined\hypertarget\hypertarget{sourcepointer-4}{}\fi
\noindent\textit{Source reference.} Original lemma 07SK: sources/stacks-spaces-more-morphisms/pinned-original-source.tex, source lines 9365--9365. Cited by the retained statement or proof but not part of the collected spans; supplied only as a navigation pointer. The argument does not depend on its content.\par\medskip
\ifdefined\hypertarget\hypertarget{sourcepointer-5}{}\fi
\noindent\textit{Source reference.} Original lemma 084T: sources/stacks-spaces-more-morphisms/pinned-original-source.tex, source lines 9365--9365. Cited by the retained statement or proof but not part of the collected spans; supplied only as a navigation pointer. The argument does not depend on its content.\par\medskip
\ifdefined\hypertarget\hypertarget{sourcepointer-6}{}\fi
\noindent\textit{Source reference.} Original lemma 0851: sources/stacks-spaces-more-morphisms/pinned-original-source.tex, source lines 9365--9365. Cited by the retained statement or proof but not part of the collected spans; supplied only as a navigation pointer. The argument does not depend on its content.\par\medskip
\ifdefined\hypertarget\hypertarget{sourcepointer-7}{}\fi
\noindent\textit{Source reference.} Original lemma 086V: sources/stacks-spaces-more-morphisms/pinned-original-source.tex, source lines 9365--9365. Cited by the retained statement or proof but not part of the collected spans; supplied only as a navigation pointer. The argument does not depend on its content.\par\medskip
\ifdefined\hypertarget\hypertarget{sourcepointer-8}{}\fi
\noindent\textit{Source reference.} Original lemma 07SU: sources/stacks-spaces-more-morphisms/pinned-original-source.tex, source lines 9365--9365. Cited by the retained statement or proof but not part of the collected spans; supplied only as a navigation pointer. The argument does not depend on its content.\par\medskip
\ifdefined\hypertarget\hypertarget{sourcepointer-9}{}\fi
\noindent\textit{Source reference.} Original lemma 082Z: sources/stacks-spaces-more-morphisms/pinned-original-source.tex, source lines 9365--9365. Cited by the retained statement or proof but not part of the collected spans; supplied only as a navigation pointer. The argument does not depend on its content.\par\medskip
\begin{lemma}[Chow's lemma]
\label{lemma-chow-noetherian-separated}
\begin{reference}
\cite[IV Theorem 3.1]{Kn}
\end{reference}
Let $S$ be a scheme. Let $f : X \to Y$ be a morphism of algebraic spaces
over $S$. Assume $f$ separated of finite type, and $Y$ separated and
Noetherian. Then there exists a commutative diagram
$$
\xymatrix{
X \ar[rd] & X' \ar[l] \ar[d] \ar[r] & \mathbf{P}^n_Y \ar[ld] \\
& Y
}
$$
where $X' \to X$ is a $U$-admissible blowup for some dense open
$U \subset X$ and the morphism $X' \to \mathbf{P}^n_Y$ is an immersion.
\end{lemma}

\begin{proof}
In this first paragraph of the proof we reduce the lemma to the case
where $Y$ is of finite type over $\Spec(\mathbf{Z})$.
We may and do replace the base scheme $S$ by $\Spec(\mathbf{Z})$.
We can write $Y = \lim Y_i$ as a directed limit of separated
algebraic spaces of finite type over $\Spec(\mathbf{Z})$, see
Limits of Spaces, Proposition \ref{spaces-limits-proposition-approximate} and
Lemma \ref{spaces-limits-lemma-descend-separated}.
For all $i$ sufficiently large we can find a separated finite type morphism
$X_i \to Y_i$ such that $X = Y \times_{Y_i} X_i$, see
Limits of Spaces, Lemmas
\ref{spaces-limits-lemma-descend-finite-presentation} and
\ref{spaces-limits-lemma-descend-separated-morphism}.
Let $\eta_1, \ldots, \eta_n$ be the generic points of the irreducible
components of $|X|$ ($X$ is Noetherian as a finite type separated
algebraic space over the Noetherian algebraic space $Y$ and therefore
$|X|$ is a Noetherian topological space).
By Limits of Spaces, Lemma \ref{spaces-limits-lemma-topology-limit}
we find that the images of $\eta_1, \ldots, \eta_n$ in $|X_i|$
are distinct for $i$ large enough. We may replace
$X_i$ by the scheme theoretic image of the (quasi-compact, in fact affine)
morphism $X \to X_i$.
After this replacement we see that the images
of $\eta_1, \ldots, \eta_n$ in $|X_i|$ are the generic points of the
irreducible components of $|X_i|$, see
Morphisms of Spaces, Lemma
\ref{spaces-morphisms-lemma-quasi-compact-scheme-theoretic-image}.
Having said this, suppose we can find a diagram
$$
\xymatrix{
X_i \ar[rd] & X_i' \ar[l] \ar[d] \ar[r] & \mathbf{P}^n_{Y_i} \ar[ld] \\
& Y
}
$$
where $X_i' \to X_i$ is a $U_i$-admissible blowup for some dense open
$U_i \subset X_i$ and the morphism $X_i' \to \mathbf{P}^n_{Y_i}$
is an immersion. Then the strict transform $X' \to X$ of $X$ relative
to $X_i' \to X_i$ is a $U$-admissible blowing up where $U \subset X$
is the inverse image of $U_i$ in $X$. Because of our carefully chosen
index $i$ it follows that $\eta_1, \ldots, \eta_n \in |U|$ and
$U \subset X$ is dense. Moreover, $X' \to \mathbf{P}^n_Y$ is an
immersion as $X'$ is closed in
$X_i' \times_{X_i} X = X_i' \times_{Y_i} Y$
which comes with an immersion into $\mathbf{P}^n_Y$. Thus we have reduced
to the situation of the following paragraph.

\medskip\noindent
Assume that $Y$ is separated of finite type over $\Spec(\mathbf{Z})$.
Then $X \to \Spec(\mathbf{Z})$ is separated of finite type as well.
We apply Lemma \ref{lemma-chow-noetherian} to $X \to \Spec(\mathbf{Z})$
to find a dense open subspace $U \subset X$ and a commutative diagram
$$
\xymatrix{
& U \ar[ld] \ar[d] \ar[rd] \ar[rrd] \\
X \ar[rd] & X' \ar[l] \ar[d] \ar[r] & Z' \ar[ld] \ar[r] & Z \ar[ld] \\
& \Spec(\mathbf{Z}) & \mathbf{P}^n_\mathbf{Z} \ar[l]
}
$$
with all the properties listed in the lemma.
Note that $Z$ has an ample invertible sheaf, namely
$\mathcal{O}_{\mathbf{P}^n}(1)|_Z$. Hence $Z' \to Z$
is a H-projective morphism by Morphisms, Lemma
\ref{morphisms-lemma-projective-over-quasi-projective-is-H-projective}.
It follows that $Z' \to \Spec(\mathbf{Z})$ is H-projective
by Morphisms, Lemma \ref{morphisms-lemma-H-projective-composition}.
Thus there exists a closed immersion
$Z' \to \mathbf{P}^m_{\Spec(\mathbf{Z})}$ for some $m \geq 0$.
It follows that the diagonal morphism
$$
X' \to Y \times \mathbf{P}^m_\mathbf{Z} = \mathbf{P}^m_Y
$$
is an immersion (because the composition with the projection
to $\mathbf{P}^m_\mathbf{Z}$ is an immersion) and we win.
\end{proof}

\begin{thebibliography}{1}
\bibitem{Kn}
Knutson, Donald.
\newblock {\em Algebraic spaces}.
\newblock Volume 203 of Lecture Notes in Mathematics. Springer-Verlag, 1971.
\end{thebibliography}
\end{document}
